\chapter{Sử dụng công cụ AI}
\label{app:ai}

\section{Phạm vi hỗ trợ}
OpenAI Codex được sử dụng để lập kế hoạch, chia việc, sinh và chỉnh sửa mã Python/Pygame, xử lý lỗi và soạn tài liệu.

Đối với mô-đun biểu diễn tri thức và mã hóa CNF (\texttt{griductive/logic/encoder.py}), Google Gemini được sử dụng ở vai trò trợ lý trao đổi lý thuyết và sinh mã kiểm thử tự động: giải thích cơ sở biến đổi công thức logic mệnh đề, rà soát hợp đồng giao diện \texttt{encode(state: PublicState) -> CNFResult}, và tạo các bộ kiểm thử \texttt{pytest} (\texttt{test\_easy\_encoder.py}, \texttt{test\_hard\_encoder.py}, \texttt{test\_superhard\_encoder.py}, \texttt{test\_integration.py}).

Đối với mô-đun bộ giải SAT bằng DPLL (\texttt{griductive/sat/dpll.py}), trợ lý AI tích hợp trong Antigravity IDE (Google Gemini) được sử dụng để giải thích lý thuyết từng bước về lan truyền đơn vị, cô lập giả định và quay lui, rà soát lỗi trong quá trình tự cài đặt mã nguồn, đề xuất và giải thích chiến lược rẽ nhánh mở rộng Dynamic Two-Sided Jeroslow--Wang (\texttt{jw}), sinh bộ kiểm thử tự động \texttt{tests/sat/test\_dpll.py} (gồm kiểm chứng bằng bảng chân trị $2^n$), và biên soạn tài liệu kỹ thuật.

\section{Một số yêu cầu tiêu biểu}
Các yêu cầu dưới đây được tóm lược từ lịch sử làm việc (không phải trích dẫn nguyên văn):
\begin{itemize}
    \item Tạo cấu trúc thư mục trò chơi và tệp \texttt{TODO.md}, chia nhiệm vụ cho năm người.
    \item Xây dựng engine và giao diện trò chơi, tách trạng thái công khai để đồng đội có thể tích hợp bộ phân tích.
    \item Sử dụng ảnh nhân vật đã cung cấp; bổ sung hoạt ảnh mở thẻ, bảng manh mối, chọn màn, điều khiển bàn phím và nhạc nền.
    \item Viết chương trình sinh màn chơi và kiểm tra màn chơi có hợp lệ với luật chơi và có lời giải không.
    \item Giải thích bài toán Griductive, cách biểu diễn CNF.
    \item Sinh các bộ kiểm thử vét cạn bảng chân trị và gán biến phụ cho các loại manh mối.
    \item Sinh bộ kiểm thử tích hợp trên toàn bộ $13$ màn chơi thực tế để đối chiếu tính tái lập, tính duy nhất nghiệm và độ tương đương suy diễn với \texttt{ReferenceAnalyzer}.
    \item Giải thích thuật toán DPLL, rà soát khung cài đặt \texttt{dpll.py} và các lỗi logic ban đầu trong lan truyền đơn vị, cô lập giả định và xử lý \texttt{timeout}.
    \item Sinh bộ kiểm thử đơn vị và đối chiếu bảng chân trị cho bộ giải DPLL; đề xuất, giải thích và đối chứng chiến lược chọn biến \texttt{jw} với \texttt{smallest}.
\end{itemize}

\section{Trách nhiệm và nguồn tài nguyên}
Luật chơi tham khảo từ Griductive~\cite{griductive} và đề bài. Ảnh, phông chữ và âm nhạc được cung cấp từ bên ngoài, không do AI tạo; nguồn được ghi trong \path{assets/ATTRIBUTION.md}~\cite{assetRecord}.

Đối với mô-đun mã hóa CNF, toàn bộ mã nguồn sản phẩm trong \texttt{griductive/logic/encoder.py} do Người~2 tự thiết kế và lập trình hoàn toàn, không đưa mã nguồn sản phẩm cho AI hay dùng AI sinh mã cài đặt. Các tệp kiểm thử do AI hỗ trợ tạo ra đều được rà soát và chạy nghiệm thu cục bộ, đạt kết quả $189/189$ ca kiểm thử thành công.

Đối với mô-đun bộ giải DPLL, Người~3 tự thiết kế, lập trình và gỡ lỗi toàn bộ phần lõi trong \texttt{griductive/sat/dpll.py} (cấu trúc \texttt{SatResult}, \texttt{\_inspect\_clause}, \texttt{\_unit\_propagate}, \texttt{\_dpll} với chiến lược cơ sở \texttt{smallest} và cơ chế \texttt{timeout}), đồng thời rà soát và kiểm định phần mở rộng chiến lược \texttt{jw} cùng toàn bộ $10/10$ ca kiểm thử trong \texttt{tests/sat/test\_dpll.py}.

Đối với mô-đun Tác tử suy diễn logic (\texttt{griductive/agent/}):
\begin{itemize}
    \item \textbf{Về trách nhiệm lập trình:} Chịu trách nhiệm chính về toàn bộ thiết kế kiến trúc \texttt{adapter.py}, quy trình suy diễn phản chứng, bộ ghi vết \texttt{trace.py} và thuật toán kiểm tra tính duy nhất \texttt{uniqueness.py}. Công cụ AI được sử dụng với vai trò trợ lý thảo luận lý thuyết logic, rà soát cú pháp, sửa lỗi logic và hỗ trợ gỡ lỗi thụt dòng. Mã nguồn cốt lõi được thành viên tự cài đặt, kiểm chứng và chạy nghiệm thu thành công trên toàn bộ các màn chơi của đề tài.
    \item \textbf{Về tài liệu tham khảo:} Lý thuyết về suy diễn phản chứng trong logic mệnh đề và nguyên lý vòng lặp suy luận được tham khảo từ sách giáo trình \textit{Artificial Intelligence: A Modern Approach} (AIMA) của Russell \& Norvig~\cite{russell2020aima}.
\end{itemize}
